\begin{textsample}{POS Dim 3 – ro – Score 21.00 – ro/ro\_inf\_28.txt}  \label{ex:f3_pos_125}
Traços, \textbf{sons}, cores e formas Este campo oportuniza experiências voltadas à expressividade das \textbf{crianças} no âmbito das artes visuais e da música, diálogo com a dança e teatro ( discutidos com mais detalhes no campo corpo, \textbf{gestos} e movimentos ), bem como com a literatura ( reportada no campo escuta, fala, pensamento e imaginação ).
Nesses processos, é preciso garantir a todas as \textbf{crianças} oportunidades para viver o prazer de pesquisar, experimentar, criar com cores, \textbf{sons}, traços e formas marcantes, que traduz a visualidade e a sonoridade presentes nestas expressões artísticas.
As \textbf{crianças} precisam ter oportunidade de conviver com diferentes manifestações artísticas, culturais e científicas, locais, regionais e universais, no cotidiano da instituição, possibilitando a apreciação, a reflexão e a criação.
Na apreciação desenvolve a \textbf{escuta} e olhar sensível, o prazer de apreciar obras musicais variadas, de diversos gêneros, estilos, épocas e culturas e das artes visuais brasileira e de outros povos, a \textbf{criança} amplia o seu repertório, analisa e estabelece relação com o que já sabe, construindo novos conhecimentos e valorizando o patrimônio cultural da humanidade.
A reflexão é considerada o complemento da apreciação, é o refletir sobre o que viu, ouviu e até mesmo sentiu, contextualização histórico-social e cultural dos produtores de artes, como forma de conhecer e representar seu povo, sua história, suas vidas em determinado período histórico da sociedade.
A criação é centrada na experimentação, na \textbf{imitação}, tendo como produtos musicais a interpretação, improvisação e a composição.
É a oportunidade que a \textbf{criança} tem de conhecer e escolher diferentes cores, \textbf{gestos}, formatos, \textbf{sons}, movimentos corporais, \textbf{texturas}, tamanhos, entre outros, por meio de experiências diversificadas, vivenciando diversas formas de expressão e linguagens artísticas.
A arte é importante meio da expressividade humana, externando sentimentos, ideias, imaginação, emoções, pensamentos, concretizando -se através das linguagens.
Com base nessas experiências, as \textbf{crianças} criam suas próprias produções artísticas ou culturais, exercitando a autoria ( coletiva e individual ) com \textbf{sons}, traços, \textbf{gestos}, danças, \textbf{mímicas}, encenações, canções, desenhos, modelagens, manipulação de diversos materiais e recursos tecnológicos.
Assim, os profissionais que atuam na Educação \textbf{Infantil} cada vez mais precisam considerar o ambiente como espaço de vivências e explorações, zona de múltiplos recursos e possibilidades para a \textbf{criança} interagir com objetos ( linhas, material não estruturado, formas, cores, aromas, \textbf{texturas}, etc. ), bem como a arte com o corpo e \textbf{sons} \textbf{emitidos} pelo próprio corpo, vivenciar sentimentos provocados pelas situações e significá -las de um modo que faça sentido para si, ampliando seu próprio repertório.
O cotidiano das \textbf{crianças} deve ser permeado de condições para sentirem a \textbf{textura} da terra ou \textbf{areia}, criar misturas, colecionar coisas, modelar com argila, criar \textbf{tintas}, explorar formas coloridas, \textbf{texturas}, sabores, \textbf{sons} e também silêncios, em um espaço acolhedor, \textbf{cheio} de visualidades e sonoridades, em espaços amplos, externos para experimentar, criar e vivenciar com o que a natureza oferece ; espaços organizados como ateliers e instalações sonoras/parques sonoros para explorações de diferentes \textbf{sons}, promovendo o desenvolvimento da expressividade e da criatividade \textbf{infantil}.
Torna -se fundamental reconhecer que cada \textbf{criança} vivencia de maneira diferente as diversas provocações feitas pelos materiais oferecidos a elas, como as significações em descascar uma fruta no momento do lanche para que as \textbf{crianças} possam percebê -la e conhecê -la no seu formato original, seu cheiro, cor de sua casca, além de seu sabor ( por exemplo, mexerica ), e outras frutas em fatias, são momentos tão ricos e diversificados quanto observar uma obra de arte.
As experiências \textbf{infantis} devem ser valorizadas na subjetividade artística.
As produções não podem ser comparadas, as \textbf{crianças} precisam ser encorajadas a se expressar pelas linguagens da arte, rompendo, com \textbf{pinturas} de desenhos prontos e estereotipados, passando a se ver como protagonistas em suas criações.
Ao oportunizar experiências musicais é importante que docentes apresentem para as \textbf{crianças} canções, \textbf{brincadeiras} cantadas, parlendas, rimas e outros jogos musicais, cantando em diferentes situações ou promovendo momentos em que todos cantam, acompanhados ou não por objetos e instrumentos musicais.
Para tanto, são valiosos os momentos em que as \textbf{crianças} observam \textbf{adultos} e outras \textbf{crianças} em apresentações e/ou improvisações musicais e festas populares.
Proporcionar de forma sistemática um repertório musical e objetos sonoros e/ou instrumentos musicais, favorecendo a exploração deles pelas \textbf{crianças} na intenção de identificar qualidades como : duração ( \textbf{sons} curtos ou longos ), altura ( \textbf{sons} graves ou agudos ), intensidade ( \textbf{sons} fracos ou fortes ) ou timbre ( que qualifica os \textbf{sons} a partir da fonte que os origina ), ampliando de forma significativa seu repertório de referências sonoras, e seus modos de escutar e produzir música.
Desenvolvendo, dessa forma, a capacidade de \textbf{brincar} com a música, de \textbf{imitar} e \textbf{inventar} movimentos que manifestam a expressividade corpórea, apreciando, produzindo e reproduzindo criações artísticas diversas.
Compreendendo a arte musical como manifestação cultural dos povos em sua construção histórica, política e social, \textbf{demonstrando} -se sensíveis à música como forma de expressão humana.
Exploração de materiais como grafite, \textbf{tintas} caseiras, sagus, guache, aquarela, giz, etc. amplia o conhecimento das \textbf{crianças} sobre o desenho e a \textbf{pintura}, conforme utilizam vários instrumentos para desenhar e pintar ( pincéis, \textbf{lápis}, carvão, canetas hidrográficas, rolinhos, escovas, esponjas, gravetos, galhinhos, barbantes, entre outros ), em diferentes planos ( horizontal e vertical ), dimensões ( bidimensional e tridimensional ) e superfícies/suportes ( telas, papéis, \textbf{paredes}, objetos de diferentes tamanhos, formatos e \textbf{texturas}, inclusive o próprio corpo ), descobrindo linhas, formas, cores, volumes, planos e utilizando -os para expressar emoções ou representar objetos e situações vividos ou imaginados.
Portanto, é extremamente importante que as \textbf{crianças} tenham a oportunidade de ampliar as capacidades autônomas de desenhar, fazendo representações de histórias, cenas ouvidas ou vivenciadas, modelar, pintar, realizar colagens etc., quando de seu interesse, predispondo -se a apreciar as próprias construções artísticas e de outros criadores ; manifestando sensibilidade à produção artística das mais diversas categorias, \textbf{demonstrando} compreensão de que a arte visual é também uma fonte de expressão do sentimento do ser humano o decorrer do seu desenvolvimento histórico e cultural.

% matched lemmas: adulto, areia, brincadeira, brincar, cheio, criança, demonstrar, emitir, escuta, gesto, imitar, imitação, infantil, inventar, lápis, mímico, parede, pintura, som, textura, tinta
\end{textsample}
